\section{The causal-support problem}

Here, \emph{causal support} is an operational training criterion, not a claim of formal
causal identification. A target action is causally supported when the variables that
justify the correction are present in, or can be obtained by an allowed action from, the
student's deployment state. Figure~\ref{fig:causal-support-gap} contrasts a supported
correction with an impossible hidden-information target.

\begin{figure}[ht]
  \centering
  \begin{minipage}[t]{0.47\textwidth}
    \centering
    \textbf{Supported correction}\\[1mm]
    \begin{tikzpicture}[node distance=5mm]
      \node[supportedbox,text width=50mm] (state)
        {Visible state $x$\\tool list: Search, ReadFile\\result: RunSQL not found};
      \node[supportedbox,text width=50mm,below=of state] (hint)
        {Grounded hint $h$\\use Search or ReadFile};
      \node[actionbox,text width=50mm,below=of hint] (action)
        {Corrected action\\call an available tool};
      \draw[proposalarrow] (state) -- node[right,font=\scriptsize]{infer} (hint);
      \draw[proposalarrow] (hint) -- node[right,font=\scriptsize]{condition} (action);
    \end{tikzpicture}

    \smallskip
    \small The hint is recoverable from the same evidence available at deployment.
  \end{minipage}
  \hfill
  \begin{minipage}[t]{0.49\textwidth}
    \centering
    \textbf{Unsupported privileged target}\\[1mm]
    \begin{tikzpicture}
      \node[proposalbox,text width=45mm] (student) at (0,0)
        {Student observes only $x$:\\request failed; no error body};
      \node[riskbox,text width=45mm] (conflict) at (0,-1.35)
        {One $P_\theta(\cdot\mid x)$ receives\\two incompatible KL targets};
      \node[actionbox,text width=24mm] (aa) at (-1.65,-2.85)
        {Hidden $z=A$\\refresh token};
      \node[actionbox,text width=24mm] (ab) at (1.65,-2.85)
        {Hidden $z=B$\\change URL};
      \draw[riskarrow] (student) -- (conflict);
      \draw[riskarrow] (conflict) -- (aa);
      \draw[riskarrow] (conflict) -- (ab);
    \end{tikzpicture}

    \smallskip
    \small The teacher sees $z$; the student cannot know which incompatible action is justified.
  \end{minipage}
  \caption{\textbf{Why privileged conditioning can create unsupported supervision.}
  OPSD is well supported when the hint is inferable from the student's state (left).
  When a hidden variable produces incompatible teacher targets for the same observed
  state (right), pulling the student toward either target cannot teach the missing
  dependency.}
  \label{fig:causal-support-gap}
\end{figure}

The failure is clearest under counterfactual variation. Suppose two training examples have
the same visible prefix $x$ but different hidden evidence $z$. The teacher sees $z$ and
therefore prefers different continuations $Q(\cdot\mid x,z)$. The student implements only
$P_\theta(\cdot\mid x)$. Minimizing divergence to both teachers cannot recover the
conditional policy; it can only fit a compromise weighted by the training frequency of
the hidden cases. On an imbalanced dataset, that compromise may look capable while merely
repeating the majority answer. Under a distribution shift, it appears as hallucination.

This problem has been observed in adjacent form: privileged self-distillation can leak
answer information and destabilize long training when the teacher signal is not anchored
by environment-derived correctness \citep{yang2026selfdistilledrlvr}. The appropriate
response is not to discard dense self-distillation. It is to separate three cases:

\begin{enumerate}
  \item \textbf{Restated evidence:} the hint names information already in $x$. OPSD and
  hinted SFT are both plausible.
  \item \textbf{Obtainable evidence:} the information is not yet in $x$ but can be acquired
  with an allowed tool or observation. The model should learn to acquire it before
  generating the hint.
  \item \textbf{Unavailable evidence:} the information exists only in a private answer,
  hidden test, or future observation. Neither method should claim autonomous causal
  competence; the system must expose the evidence, defer, or abstain.
\end{enumerate}
